% =====================================================================
\chapter{Первое включение и настройки пульта}
\label{ch:radio}
% =====================================================================
\chapterintro
  {что за предупреждения появляются при включении, как откалибровать стики, выбрать
   режим стиков и порядок каналов, настроить звук, голос, подсветку и контроль батареи}
  {пульт с заряженными аккумуляторами}

\section{Предупреждения при включении}

\begin{simple}
Пульт при включении проверяет, не опасно ли сейчас начинать. Если стик газа поднят или
переключатели стоят «не так», он не пустит тебя дальше, пока ты не исправишь. Это как
ремень безопасности в машине: раздражает, но спасает.
\end{simple}

\begin{figure}[H]
\centering
\begin{tabular}{@{}c@{\hspace{6mm}}c@{}}
\lcdscreen{THROTTLE WARNING}{}{\lblank\lc{\bfseries Throttle not idle}\lblank\lc{Опусти газ вниз}\lblank\lc{\linv{ Press any key to skip }}} &
\lcdscreen{SWITCH WARNING}{}{\lblank\lc{\bfseries Switches not in}\lc{\bfseries default position}\lblank\lc{\linv{SA↓}\ \ \linv{SC-}}\lblank}
\end{tabular}
\caption{Слева: газ не внизу. Справа: переключатели \sw{SA} и \sw{SC} не в стартовом положении}
\end{figure}

\begin{description}[font=\sffamily\bfseries\color{etx}]
  \item[Throttle warning] опусти стик газа до конца вниз --- надпись исчезнет сама.
  \item[Switch warning] переведи показанные переключатели в стартовое положение
        (обычно --- все вверх).
  \item[SD card warning] версия файлов на SD-карте не совпадает с прошивкой
        (глава~\ref{ch:firmware}).
  \item[Sound off] звук выключен --- ты можешь не услышать важное предупреждение.
\end{description}

\begin{caution}
Не привыкай пропускать предупреждения кнопкой. Однажды модель с поднятым газом рванёт
с места прямо у тебя в руках.
\end{caution}

\section{Калибровка стиков}

\begin{simple}
Пульт должен точно знать, где у каждого стика центр и где края. Этому его учит
калибровка. Делай её на новом пульте, после прошивки и если стики «врут».
\end{simple}

\begin{figure}[H]
\centering
\begin{tabular}{@{}ccc@{}}
\begin{tikzpicture}[scale=0.75]
  \pic[transform shape] at (0,0) {stickidle};
  \pic[transform shape] at (2.5,0) {stickidle};
  \node[font=\sffamily\scriptsize,align=center] at (1.25,-1.5) {отпусти стики,\\крутилки --- в середину};
\end{tikzpicture} &
\begin{tikzpicture}[scale=0.75]
  \foreach \x in {0,2.5}{
    \fill[black!8,draw=black!40,rounded corners=1.2mm] (\x-1,-1) rectangle (\x+1,1);
    \draw[motion,-{Latex[length=2mm]}] (\x+0.8,0) arc (0:340:0.8);
    \shade[ball color=black!55] (\x+0.8,0) circle (0.2);}
  \node[font=\sffamily\scriptsize,align=center] at (1.25,-1.5) {обведи стики по кругу\\до упора 2--3 раза};
\end{tikzpicture} &
\begin{tikzpicture}[scale=0.75]
  \pic[transform shape] at (0,0.1) {knob=-135};
  \pic[transform shape] at (2.2,0.1) {knob=135};
  \draw[motion,-{Latex[length=2mm]}] (0.5,0.9) -- (1.7,0.9);
  \node[font=\sffamily\scriptsize,align=center] at (1.1,-1.5) {прокрути крутилки\\от края до края};
\end{tikzpicture}\\
\small\bfseries шаг 1 & \small\bfseries шаг 2 & \small\bfseries шаг 3
\end{tabular}
\caption{Калибровка в три шага}
\end{figure}

\begin{enumerate}
  \item \mpath{SYS\sep Hardware}, пункт \menu{Calibration} (в некоторых версиях --- отдельная
        страница \menu{Calibration}). Нажми \btn{ENTER}.
  \item Отпусти стики в центр, крутилки поставь в середину → \btn{ENTER}.
  \item Обведи стики по кругу до упора, прокрути крутилки → \btn{ENTER}. Готово.
\end{enumerate}

\begin{tip}
Проверь на главном экране (страница с каналами): в центре стики должны давать 0, по краям ---
ровно $-100$ и $+100$. Если газ внизу показывает, например, $-96$, повтори калибровку.
\end{tip}

\section{Режим стиков и порядок каналов}

\mpath{SYS\sep Radio Setup} --- длинная страница общих настроек. Две самые важные:

\begin{description}[font=\sffamily\bfseries\color{etx}]
  \item[Mode] режим стиков 1--4 (глава~\ref{ch:axes}). Почти всегда \textbf{2}.
  \item[Default channel order] в каком порядке стики попадут в первые четыре канала
        \emph{у новых моделей}.
\end{description}

\begin{figure}[H]
\centering
\begin{tikzpicture}[x=13mm,y=8mm,every node/.style={font=\sffamily\small}]
  \node[anchor=east,font=\sffamily\small\bfseries] at (0.3,0) {AETR:};
  \foreach \n/\c [count=\i] in {Ail/etx,Ele/etx,Thr/accent,Rud/etx}{
    \node[draw=\c,fill=\c!12,rounded corners=2pt,minimum width=11mm,minimum height=6mm] at (\i,0) {\n};
    \node[font=\sffamily\scriptsize,text=black!60] at (\i,0.65) {CH\i};}
  \node[anchor=west,font=\sffamily\scriptsize] at (4.6,0) {ExpressLRS, Betaflight, INAV --- \textbf{выбирай это}};
  \node[anchor=east,font=\sffamily\small\bfseries] at (0.3,-1.2) {TAER:};
  \foreach \n/\c [count=\i] in {Thr/accent,Ail/etx,Ele/etx,Rud/etx}{
    \node[draw=\c,fill=\c!12,rounded corners=2pt,minimum width=11mm,minimum height=6mm] at (\i,-1.2) {\n};}
  \node[anchor=west,font=\sffamily\scriptsize] at (4.6,-1.2) {порядок Spektrum/JR};
\end{tikzpicture}
\caption{Два популярных порядка каналов. Газ (оранжевый) --- на третьем или первом месте}
\end{figure}

\begin{note}
Порядок каналов пульта и порядок, который ждёт модель (например, \menu{Channel map} в
Betaflight), \textbf{должны совпадать}. Иначе ты будешь давать газ, а коптер --- крениться.
\end{note}

\section{Батарея пульта}

\begin{center}
\lcdscreen{RADIO SETUP}{3/9}{\lr{Battery range}{6.6-8.4V}\lsel{Battery low}{\ledit{7.0V}}\lr{Battery calib}{7.9V}\lr{Inactivity alarm}{10m}\lr{Sound mode}{All}\lblank}
\end{center}

\begin{itemize}
  \item \menu{Battery range} --- от «пустой» до «полной» батареи на значке. Для двух банок
        (2×18650 или 2S LiPo): \textbf{6,6--8,4\,В}.
  \item \menu{Battery low} --- когда пульт начнёт пищать «батарея разряжена»:
        около \textbf{7,0\,В}.
  \item \menu{Battery calib} --- если показания врут, измерь напряжение мультиметром и
        выставь такое же.
\end{itemize}

\section{Звук, голос, вибрация}

\begin{itemize}
  \item \menu{Sound mode} --- что озвучивать. Режим \menu{All} включает всё, включая голос.
  \item \menu{Volume} и другие громкости --- громкость голоса, писков, музыки.
  \item \menu{Voice language} --- язык голоса. Нужна папка \code{SOUNDS/ru} (для русского)
        на SD-карте --- её устанавливает Buddy (глава~\ref{ch:firmware}).
  \item \menu{Haptic} --- вибромотор: включить, сила, длительность.
\end{itemize}

\section{Подсветка, время и прочее}

\begin{itemize}
  \item \menu{Backlight} --- когда включать подсветку экрана и на сколько секунд.
        Подсветка расходует батарею --- 10--30 секунд достаточно.
  \item \menu{Date}, \menu{Time} --- для правильных имён файлов логов.
  \item \menu{Inactivity alarm} --- напомнит о забытом включённом пульте.
  \item \menu{USB mode} --- что делать при подключении USB: \menu{Ask} (спросить),
        \menu{Joystick} (для симуляторов), \menu{Storage} (как флешка), \menu{Serial}.
        Удобнее всего оставить \menu{Ask}.
  \item \menu{Owner} --- имя владельца. Полезно в кружке, где много одинаковых пультов.
\end{itemize}

\section{Hardware: «паспорт» пульта}

\mpath{SYS\sep Hardware} описывает, какие переключатели и крутилки физически стоят в пульте.

\begin{center}
\lcdscreen{HARDWARE}{8/9}{\lr{SA}{2POS}\lr{SB}{3POS}\lr{SC}{3POS}\lsel{SD}{\ledit{Toggle}}\lr{S1}{Pot w. det}\lr{Internal RF}{CRSF}}
\end{center}

\begin{itemize}
  \item тип переключателя: \menu{2POS}, \menu{3POS}, \menu{Toggle} (без фиксации) или
        \menu{None} (нет);
  \item тип крутилки: с фиксацией в середине или без;
  \item имя: можно переименовать, например \sw{S1} → \sw{Cam};
  \item \menu{Internal RF} --- тип встроенного радиомодуля (\menu{CRSF} для ELRS,
        \menu{Multi} для «4-в-1»).
\end{itemize}

\begin{caution}
Если тип переключателя указан неправильно, среднее положение не будет работать. Если
неправильно указан тип внутреннего модуля, пульт не сможет с ним говорить. Меняй здесь
что-то, только если точно знаешь, что делаешь.
\end{caution}

\begin{tryit}
\begin{enumerate}
  \item Откалибруй стики и проверь на странице каналов: 0 в центре, $\pm100$ по краям.
  \item Установи \menu{Battery low} = 7,0\,В.
  \item Выставь \menu{Default channel order} = AETR и \menu{Mode} = 2.
  \item Включи русский голос (если он установлен) и проверь, говорит ли пульт.
\end{enumerate}
\end{tryit}

\begin{summary}
\begin{itemize}
  \item Предупреждения при включении --- для безопасности; исправляй, а не пропускай.
  \item Калибровка: центр → круги до упора → сохранить.
  \item Mode 2 и порядок каналов AETR --- стандарт для ELRS и квадрокоптеров.
  \item Батарея 2S: шкала 6,6--8,4\,В, предупреждение около 7,0\,В.
\end{itemize}
\end{summary}
